W celu zapewnienia rzetelnego punktu odniesienia oraz właściwej kontekstualizacji uzyskanych wyników przeprowadzono ewaluację wybranych modeli typu state-of-the-art (SOTA), w tym modeli dostępnych w bibliotece InsightFace oraz modelu referencyjnego.

Wybrane modele obejmowały zarówno rozwiązania o wysokiej dokładności, jak i lekkie architektury zoptymalizowane pod kątem wydajności obliczeniowej, a także model referencyjny starszej generacji.

Testowane modele:
\begin{itemize}
\item ArcFace Large (buffalo\_l) – model oparty na architekturze IResNet-50, dostępny w bibliotece InsightFace,
\item ArcFace Small (buffalo\_s) – lekka architektura oparta na MobileFaceNet, przeznaczona do zastosowań mobilnych, dostępna w bibliotece InsightFace,
\item VGGFace2 (ResNet-50) – model referencyjny starszej generacji (2018 r.), zaimplementowany z wykorzystaniem biblioteki \textit{keras-vggface}, niebędący częścią frameworka InsightFace.
\end{itemize}

Wszystkie modele poddano temu samemu protokołowi identyfikacji (galeria–sonda) oraz symulacji okluzji w postaci poziomego zasłonięcia regionu oczu o wysokości 20 pikseli. Należy jednak podkreślić, że część modeli SOTA wykorzystuje własny pipeline ekstrakcji cech: dla ArcFace (InsightFace) embeddingi są wyznaczane na podstawie pełnego obrazu z detekcją twarzy, natomiast dla VGGFace (ResNet50) stosowane jest przycięcie względem \texttt{bbox} i skalowanie do 224×224. Zostało to utrzymane zgodnie z implementacjami referencyjnymi w repozytorium. Szczegółowy opis protokołu ewaluacji przedstawiono w dalszej części dokumentu.